%% ============================================================
%%  Honours FINAL PAPER SCAFFOLD — RQ3
%%  Feudal Networks vs flat PPO: recovery after opponent shifts in Go
%%
%%  ACM acmart, thesis config (sigconf two-column, nonacm = no ACM footers).
%%  Compile on Overleaf with pdfLaTeX.
%%
%%  SCAFFOLD RULES followed here:
%%   - Abstract = COMPLETE first draft.
%%   - All other sections = PLANNED: short bullets saying WHAT will be written
%%     (not the prose), with [TBD] where numbers/results are not yet final.
%% ============================================================
\documentclass[sigconf,nonacm]{acmart}
\settopmatter{printacmref=false}
\usepackage{booktabs}
\usepackage{graphicx}

\begin{document}

\title{Does a feudal reinforcement learning improve recovery time following
       abrupt opponent strategy shifts, compared to flat RL approach (PPO)?}

\author{Brad Scott}
\affiliation{%
  \institution{University of Cape Town}
  \department{Department of Computer Science}
  \city{Cape Town}
  \country{South Africa}}
\email{sctbra008@myuct.ac.za}

%% -------------------- ABSTRACT: complete first draft --------------------
\begin{abstract}
Reinforcement learning is increasingly applied to competitive settings in which
an opponent's behaviour is not fixed. When that behaviour changes abruptly, the
resulting non-stationarity can invalidate a previously effective policy.
Hierarchical architectures are widely expected to re-adapt faster than flat ones
in such situations, yet this expectation has rarely been tested directly against
a shifting opponent. This work investigates whether a hierarchical Feudal
Network (FuN) recovers more quickly than a flat Proximal Policy Optimization
(PPO) agent following abrupt opponent strategy shifts in the game of Go. Both
agents are trained against fixed heuristic opponents in PettingZoo's
\texttt{go\_v5} on $9\times9$ and $13\times13$ boards, after which mid-training
opponent shifts of varying magnitude and frequency are introduced. Their effect
is measured by recovery time, together with the depth and overall cost of the
resulting performance drop, over multiple random seeds. Across the conditions
examined, the hierarchical agent does not recover faster than the flat baseline.
These findings suggest that the FuN hierarchy does not provide the expected
advantage under opponent-induced non-stationarity.
\end{abstract}

\maketitle

%% ==================== PLANNED SECTIONS (bullets) ====================

\section{Introduction}
Reinforcement learning (RL) is a framework in which an agent learns how to
behave by interacting with an environment. The agent observes different states,
takes actions on those states and receives a reward, gradually adjusting its
behaviour to maximise these rewards it collects over
time~\cite{sutton2018reinforcement}. Through repeated trial and error the agent
learns a policy, a function that maps each state to an action. Such a policy
works well as long as the environment stays the same, but many real settings are
not static. In competitive games in particular, an opponent does not play the
same way forever and may change its strategy as the game unfolds. When the
opponent changes, the conditions the agent originally learned against no longer
hold, so behaviour that was once effective can suddenly perform poorly. This
effect, in which the environment an agent faces shifts because the opponent's
behaviour changes, is known as
non-stationarity~\cite{hernandezleal2017survey}, and it is the central
difficulty this work addresses. An agent that cannot adjust quickly after such a
change suffers a lasting drop in performance, which motivates investigating how
the architecture of a learning agent affects how quickly it recovers once an
opponent shifts strategy.

\begin{itemize}
  \item The idea under test: a manager/worker hierarchy should recover faster.
\end{itemize}

\subsection{Research Question}
\begin{itemize}
  \item State RQ3 exactly. Describe the setting: the game of Go, against fixed
        rule-based heuristic opponents on two board sizes ($9\times9$ and
        $13\times13$).
  \item What this work adds: (i) a controlled opponent-shift recovery benchmark
        in Go; (ii) a tested FuN vs PPO comparison under identical conditions,
        across shift magnitude, frequency and board size, (iii) three recovery
        measures and what each can and cannot detect.
  \item Main finding: the hierarchy did not improve recovery.
  \item Roadmap: preview the remaining sections - Related Works, Design and
        Implementation, Results, Discussion, and Conclusions.
\end{itemize}

\section{Related Works}
% A SYNTHESISED discussion that builds to the gap — not a list of papers.
% Keep it only as long as it needs to be. Opening frame, then subsections,
% then close on the gap.
\begin{itemize}
  \item Open with: flat versus hierarchical RL, both of which must handle a
        shifting opponent - followed by a preview of subsections.
  \item Well-known RL work on Go, such as AlphaGo, as background context.
\end{itemize}

\subsection{Non-stationarity in Multi-Agent RL}
\begin{itemize}
  \item Non-stationarity / opponent-induced dynamics in multi-agent RL; why a
        shift breaks a learned policy.
\end{itemize}

\subsection{Flat and Hierarchical RL}
\begin{itemize}
  \item Flat on-policy RL: PPO - the baseline under test.
  \item Hierarchical RL and Feudal Networks: Manager--Worker, temporal
        abstraction, intrinsic sub-goal reward; note the claims that it adapts
        to a different opponent.
\end{itemize}

\subsection{Measuring Adaptation and Recovery}
\begin{itemize}
  \item Recovery-time / adaptation-speed metrics; how the recovery measure
        relates and contrasts with catastrophic-forgetting work.
\end{itemize}

% Close the section on the gap:
\begin{itemize}
  \item State the gap this paper fills: HRL vs flat recovery under an opponent
        shift, tested directly.
\end{itemize}

\section{Design and Implementation}
\begin{itemize}
  \item Open with: brief outline of the section - the environment and
        opponents, the agents, the experimental design (Phase 1 and Phase 2),
        the metrics, and the tools used.
\end{itemize}

\subsection{Environment and Opponents}
\begin{itemize}
  \item \textbf{Environment:} PettingZoo \texttt{go\_v5}; $9\times9$ (komi 5.5)
        and $13\times13$ (komi 7.5); area scoring. Note the move-cap /
        forced-termination handling.
  \item \textbf{Fixed opponents:} five heuristic bots (greedy, defensive,
        corner, edge, random), description of each style.
\end{itemize}

\subsection{Agents}
\begin{itemize}
  \item \textbf{Agents:} PPO (action-masked categorical policy, shared CNN
        backbone); FuN (matched CNN perception, Manager dilated-LSTM trained on
        extrinsic reward, Worker trained on intrinsic cosine reward). State that
        both share observation preprocessing and backbone so the comparison is
        architectural; note key implementation corrections made to FuN.
\end{itemize}

\subsection{Experimental Design}
\begin{itemize}
  \item \textbf{External reference:} GNU~Go, to reference agent strength before
        the comparison.
  \item \textbf{Phase 1:} train each agent against each opponent to a stable
        win-rate baseline; this baseline defines the recovery target.
  \item \textbf{Phase 2:} resume a Phase-1 checkpoint, then apply abrupt
        opponent shifts.
        \begin{itemize}
          \item Magnitudes, calibrated by observed disruption: LOW
                corner$\to$edge, MED corner$\to$defensive, HIGH
                greedy$\to$defensive.
          \item Frequencies: f1 single shift, f2 periodic, f3 frequent.
          \item Design table: $3$ magnitudes $\times$ $3$ frequencies $\times$
                $2$ agents $\times$ $2$ boards $\times$ $3$ seeds.
        \end{itemize}
  \item \textbf{Table of experiments:} summarise the runs in a table - columns
        for condition, what is measured, and purpose.
  \item \textbf{Parameters:} exact shift schedules and step budgets.
\end{itemize}

\subsection{Metrics}
\begin{itemize}
  \item \textbf{Metrics used:} Rolling win rate (window $100$ games); recovery
        time; dip depth; adaptation cost. Description and purpose of each metric.
\end{itemize}

\subsection{Tools}
\begin{itemize}
  \item \textbf{Setup:} GPU and CPU training on the UCT cluster and CPU training
        on the CHPC cluster; metrics logged to Weights \& Biases and to per-game
        CSV files; per-seed runs kept separate then merged; state the number of
        seeds.
\end{itemize}

\section{Results}
% Factual presentation only — no interpretation (that goes in Discussion).
\begin{itemize}
  \item Open with: brief outline of the results - Phase-1 capability baselines,
        then Phase-2 recovery, presented as tables and figures.
\end{itemize}

\subsection{Phase-1 Capability}
\begin{itemize}
  \item \textbf{Phase-1 capability} - Table: win rate per agent $\times$
        opponent $\times$ board, with the GNU~Go reference [values TBD].
\end{itemize}

\subsection{Phase-2 Recovery}
\begin{itemize}
  \item \textbf{Phase-2 recovery} - Table: recovery time, dip depth and
        adaptation cost per agent $\times$ magnitude $\times$ board, as mean
        over seeds [values TBD].
  \item \textbf{Figures:}
        \begin{itemize}
          \item rolling win-rate curves with shift points marked, PPO and FuN
                overlaid per condition;
          \item grouped bar chart of recovery time by magnitude, per board;
          \item dip-depth / adaptation-cost comparison across conditions.
        \end{itemize}
\end{itemize}

\section{Discussion}
\begin{itemize}
  \item Open with: recap what was tested, then interpret the Results.
  \item Does the result answer RQ3 / reject the hypothesis that hierarchy speeds
        recovery?
  \item Analyse \emph{why} the result is negative - candidate mechanisms.
  \item How recovery depends on magnitude, frequency and board size, and what
        the two boards reveal about the metrics themselves (recovery time is
        informative on $9\times9$; dip depth / adaptation cost on $13\times13$).
  \item How the finding aligns with or contrasts the related literature (HRL's
        claimed recovery advantage vs an opponent-shift setting).
  \item \textbf{Implications:} what the negative result means for using HRL
        under opponent-induced non-stationarity, and when a hierarchy might
        still help.
  \item \textbf{Limitations:} all seeds start from the same Phase-1 model, so
        they capture only Phase-2 randomness; win rate used rather than a
        territory/score margin; boards limited to $9\times9$ and $13\times13$,
        not full-size $19\times19$ Go.
\end{itemize}

\section{Conclusions}
\begin{itemize}
  \item Restate the answer to RQ3 in one line [direction known, final numbers
        TBD]: the FuN hierarchy did not recover faster than flat PPO.
  \item Implication for hierarchical RL under opponent-induced non-stationarity.
  \item Future work: more seeds; $19\times19$ board; a territory-margin metric.
\end{itemize}

\bibliographystyle{ACM-Reference-Format}
\bibliography{references}

\end{document}
